% TOP_PROBLEM: {"id": 30006786, "title": "Parity conjecture for elliptic curves", "category_id": 1, "category": {"id": 1, "name": "number_theory", "display_name": "Number Theory", "description": "Properties of integers, prime numbers, Diophantine equations.", "slug": "number-theory", "order_index": 1, "created_at": "2026-07-31T15:26:25.670Z"}, "status": "open"}
% FIRST_POSTED: 2026-09-27
% !TeX program = lualatex
% Compile twice: lualatex open_problems_500.tex
% All references and prose are embedded; no BibTeX database or external inputs.
\documentclass[11pt,a4paper]{article}
\usepackage[margin=24mm,headheight=14pt]{geometry}
\usepackage{fontspec}
\setmainfont{Latin Modern Roman}
\setsansfont{Latin Modern Sans}
\setmonofont{Latin Modern Mono}
\usepackage{amsmath,amssymb,mathtools}
\usepackage{unicode-math}
\setmathfont{Latin Modern Math}
\usepackage{microtype}
\usepackage{enumitem,xurl,needspace}
\usepackage[dvipsnames]{xcolor}
\usepackage{hyperref}
\hypersetup{unicode=true,colorlinks=true,linkcolor=MidnightBlue,urlcolor=MidnightBlue,
 pdftitle={500 Mathematical Problem Statements},pdfauthor={Source-annotated compilation},bookmarksnumbered=true}
\usepackage{bookmark}
\usepackage{tocloft}
\setlength{\cftsecnumwidth}{3em}
\cftsetpnumwidth{2.5em}
\cftsetrmarg{4em}
\usepackage{titlesec}
\titleformat{\section}{\Large\bfseries\raggedright}{\thesection}{0.7em}{}
\usepackage{fancyhdr}
\pagestyle{fancy}\fancyhf{}
\fancyhead[L]{\small\sffamily Mathematical problem statements}
\fancyhead[R]{\small Source edition 22}
\fancyfoot[C]{\thepage}
\setlength{\parindent}{0pt}
\setlength{\parskip}{5pt plus 1pt}
\setlength{\emergencystretch}{3em}
\setcounter{tocdepth}{1}
\setcounter{secnumdepth}{1}
\setlist[itemize]{leftmargin=3em,itemsep=5pt,topsep=4pt,parsep=0pt}
\allowdisplaybreaks
\newcommand{\N}{\mathbb{N}}
\newcommand{\Z}{\mathbb{Z}}
\newcommand{\Q}{\mathbb{Q}}
\newcommand{\R}{\mathbb{R}}
\newcommand{\C}{\mathbb{C}}
\newcommand{\F}{\mathbb{F}}
\newcommand{\A}{\mathbb{A}}
\newcommand{\PP}{\mathbb P}
\newcommand{\E}{\mathbb{E}}
\newcommand{\eps}{\varepsilon}
\newcommand{\dd}{\,\mathrm d}
\newcommand{\sref}[2]{\hyperlink{source:#1:#2}{[S#2]}}
\newcommand{\eref}[2]{\hyperlink{extra:#1:#2}{[E#2]}}
% Turn the next switch off for a shorter reading copy. The quotations preserve
% every exactTarget field, including qualifications omitted from a short summary.
\newif\ifshowcatalogue
\showcataloguetrue
\newcommand{\cataloguescope}[1]{\ifshowcatalogue
 \paragraph{Exact scope in the supplied catalogue.}
 \begin{quote}\small #1\end{quote}\fi}

% Compile twice with: lualatex 30006786.tex
\numberwithin{equation}{section}
\hypersetup{pdftitle={Research attempt -- problem 30006786},pdfauthor={Research notebook; original compilation retained}}
\fancyhead[L]{\small\sffamily Research notebook}
\fancyhead[R]{\small\sffamily Problem 30006786 | 27 September 2026}
\begin{document}
\setcounter{section}{0}
% ================================================================
% problemId: 30006786
% Canonical problem ID: 30006786
% exactTarget SHA-256: 7530843b83339e097ed6661b5a9cf2a6ee2152588c7135a0af4b7e9de37d2fe4
\clearpage
\section*{\texorpdfstring{Parity conjecture for elliptic curves}{Parity conjecture for elliptic curves}}\label{30006786}
\subsection{Definitions and mathematical statement}
By the Mordell--Weil rank of $E/K$ mean the integer $r$ in
$E(K)\cong E(K)_{\mathrm{tors}}\oplus\Z^r$. The global root number is
$w(E/K)=\prod_v w(E/K_v)\in\{\pm1\}$, the product of normalized local epsilon-factor signs; almost every finite-place factor is one and each archimedean factor is $-1$. The target is
\[
 w(E/K)=(-1)^r
\]
for every elliptic curve over every number field. Root numbers are defined locally without presuming a global analytic continuation theorem. The claim concerns the algebraic rank, not merely parity of a conjectural analytic order of vanishing or a Selmer-group rank.
\par\smallskip\textit{Formulation and concept references:} \sref{30006786}{1}.
\cataloguescope{For every number field K and every elliptic curve E/K, is w(E/K)=(-1)\textasciicircum{}rank\_Z E(K), where rank\_Z E(K) is the Mordell-Weil rank and w(E/K)=product\_v w(E/K\_v) is the global root number? The local factors are the usual Weil-Deligne root numbers at finite places, equal to1 at all but finitely many places, and -1 at every archimedean place; their definition does not assume analytic continuation of L(E/K,s).}
\subsection{Short English statement}
Does the product of an elliptic curve's local root-number signs always determine whether its rational-point rank over the number field is even or odd?
% BEGIN RESEARCH ADDITION ATTEMPT-196
\subsection{Research attempt}

\paragraph{Current frontier.}
The formulation source proves $p$-parity when $E$ has a $p$-isogeny, including an odd-degree Galois extension variant; its Corollary~1.5 adds finiteness of the $p$-primary Tate--Shafarevich group to deduce Mordell--Weil parity \sref{30006786}{1}, Theorems~1.4--1.4$'$ and Corollary~1.5. It also proves the full parity assertion when complex multiplication is defined over the ground field. Those precise statements must not be converted into unconditional Mordell--Weil parity for every $E/K$. The attempt targets the divisible Tate--Shafarevich contribution, which a superficial use of the Cassels--Tate pairing can miss.

\paragraph{Attack route.}
One route proves finiteness of the Tate--Shafarevich group, which is stronger than needed. Another proves just evenness of its divisible $p$-corank. A third combines quadratic extension and twist identities to propagate signs, but a product of two unknown parities does not determine either individually. We pursue the weaker evenness requirement and test a proposed pairing proof. The result is an exact reduction on the known $p$-parity locus and an explicit abstract model refuting the pairing shortcut.

\paragraph{Attempt.}
Fix a prime $p$. Write
\[
 r=\operatorname{rank}E(K),\quad s_p=\operatorname{corank}_{\Z_p}
       \operatorname{Sel}_{p^\infty}(E/K),\quad
 d_p=\operatorname{corank}_{\Z_p}\operatorname{Sha}(E/K)[p^\infty].
\]
The Kummer exact sequence is
\begin{equation}\label{eq196:selmer}
 0\longrightarrow E(K)\otimes\Q_p/\Z_p
 \longrightarrow\operatorname{Sel}_{p^\infty}(E/K)
 \longrightarrow\operatorname{Sha}(E/K)[p^\infty]\longrightarrow0.
\end{equation}
These cofinitely generated groups have additive corank, hence $s_p=r+d_p$; this is the source's distinction between the two ranks \sref{30006786}{1}, equations (4)--(5).

On any locus where the established $p$-parity theorem gives $w(E/K)=(-1)^{s_p}$,
\begin{equation}\label{eq196:defect}
 w(E/K)(-1)^r=(-1)^{d_p}.
\end{equation}
Thus full parity on that locus is equivalent to evenness of $d_p$, not to finiteness of the whole group. In particular $d_p=0$ is sufficient but stronger. No equality between algebraic and analytic rank has been used, so the root number retains the local definition in the target.

The proposed proof is to use the Cassels--Tate pairing. Its radical on the relevant primary part is the maximal divisible subgroup; on the finite quotient it gives the expected alternating duality for an elliptic curve. The source explicitly separates that finite quotient from the divisible contribution \sref{30006786}{1}, discussion preceding Theorem~1.10. It therefore proves an evenness statement for the finite paired part, not automatically for $d_p$.

\textbf{Abstract countermodel to the shortcut.} For any prime $p$ and integer $a\ge1$, let
\[
 A=(\Q_p/\Z_p)\oplus(\Z/p^a\Z)^2.
\]
Define a bilinear pairing into $\Q/\Z$ by
\[
 \langle(t,x,y),(t',x',y')\rangle=\frac{xy'-yx'}{p^a}\pmod\Z.
\]
It is alternating, its radical is exactly the first, divisible summand, and its quotient pairing is perfect. The finite quotient has square order $p^{2a}$ and even $p$-rank, while the divisible corank is one. Every abstract property just cited is compatible with odd $d_p$. This is not claimed to be the Tate--Shafarevich group of an elliptic curve; it disproves only the logical inference from those pairing properties.

What additional structure would suffice? Let
\[
 D_p=\operatorname{Sha}(E/K)[p^\infty]_{\rm div},\qquad
 V_p(D_p)=\left(\varprojlim_m D_p[p^m]\right)\otimes_{\Z_p}\Q_p,
\]
where transition maps are multiplication by $p$. Since
$D_p\simeq(\Q_p/\Z_p)^{d_p}$, the vector space $V_p(D_p)$ has dimension $d_p$. A nondegenerate alternating $\Q_p$-bilinear form on it would force $d_p$ even. For odd $d_p$, its matrix $B$ would satisfy
\[
 \det B=\det B^{\mathsf T}=\det(-B)=(-1)^{d_p}\det B=-\det B,
\]
so $\det B=0$ in the characteristic-zero field $\Q_p$, including $p=2$. This is a sufficient missing pairing lemma. The existing Cassels--Tate pairing cannot supply it by restriction: it vanishes identically on $D_p$. A new derived or height-theoretic construction would need both a definition on this vector space and a proof of nondegeneracy.

Even that hypothetical lemma would settle the general target only after $p$-parity is available for a suitable prime for every $E/K$. We have not established that global coverage. Within the $p$-isogeny cases covered by the source, however, the single even-corank statement would suffice by \eqref{eq196:defect}.

\paragraph{Stress test.}
Finite-level groups can have even-dimensional nondegenerate quotients while their compatible divisible kernels have odd corank; the countermodel shows this without a limiting ambiguity. Therefore taking an inverse limit of quotient pairings and declaring it nondegenerate is unjustified. Nor does a product of local root numbers supply a pairing on $D_p$ by formal analogy. For curves with ground-field complex multiplication the source gives additional arithmetic structure and actual parity; that special input cannot be extended to a general elliptic curve by ignoring endomorphism hypotheses. The finite alternating-matrix checks in the script test the countermodel, not existence of a curve with odd $d_p$.

\paragraph{Outcome.}
\textbf{Outcome: Conditional reduction.}
On the established $p$-parity locus, the exact remaining defect is $(-1)^{d_p}$. A nondegenerate alternating form on the divisible Tate module would remove it, but the usual pairing does not provide such a form, and a concrete abstract example explains why. No finiteness of $\operatorname{Sha}$, general even-corank theorem, or unconditional parity theorem beyond the cited cases is proved.

% END RESEARCH ADDITION ATTEMPT-196
\subsection{Sources}
\begin{itemize}\raggedright
\item[\textnormal{[S1]}]\hypertarget{source:30006786:1}{} Kęstutis Česnavičius, The p-parity conjecture for elliptic curves with a p-isogeny; arXiv1207.0431v3 April8,2014; PDF internal date May22,2018.
\url{https://arxiv.org/pdf/1207.0431}.
{\small\textit{Catalogue formulation source.}}
\end{itemize}

\end{document}
